%hw-3.tex
%%%%%%%%%%%%%%%%%%%%%%%%%%%%%%%%%%%%%%%%%%%%%%%%%%%%%%%%%%%%%%%%%%%%%%
%%%%%%            Math 403 Homework \#3, Spring 2025            %%%%%%
%%%%%%                                                          %%%%%%
%%%%%%                 Instructor: Ezra Miller                  %%%%%%
%%%%%%%%%%%%%%%%%%%%%%%%%%%%%%%%%%%%%%%%%%%%%%%%%%%%%%%%%%%%%%%%%%%%%%
\documentclass[11pt]{article}

\oddsidemargin=17pt \evensidemargin=17pt
\headheight=9pt     \topmargin=26pt
\textheight=564pt   \textwidth=433.8pt
\parskip=1ex        %voffset=-6ex

\usepackage{tikz,amsmath,amssymb,graphicx,color,enumitem,graphicx}%,setspace,mathrsfs}%,hyperref
\setlist[enumerate]{parsep=1.5ex plus 4pt,topsep=0ex plus 4pt,itemsep=0ex}
  \setlist[itemize]{parsep=1.5ex plus 4pt,topsep=-1ex plus 4pt,itemsep=-1.33ex}

\leftmargini=5.5ex
\leftmarginii=3.5ex

%for \marginpar to fit optimally
\hoffset=-1.02in
\setlength\marginparwidth{2.2in}
\setlength\marginparsep{1mm}
\newcommand\red[1]{\marginpar{\linespread{.85}\sf%
		   \vspace{-1.4ex}\footnotesize{\color{red}#1}}}
\newcommand\score[1]{\marginpar{\vspace{-2ex}\color{blue}{#1/3}}\hspace{-1ex}}
\newcommand\extra[1]{\marginpar{\color{blue}{#1/1}}\hspace{-1ex}}
\newcommand\total[2]{\marginpar{\colorbox{yellow}{\huge #1/#2}}}
\newcommand\collab[1]{\marginpar{\vspace{-11ex}\colorbox{yellow}{#1/3}}\hspace{-1ex}}
\newcommand\magenta[1]{\colorbox{magenta}{\!\!#1\!\!}}
\newcommand\yellow[1]{\colorbox{yellow}{\!\!#1\!\!}}
\newcommand\green[1]{\colorbox{green}{\!\!#1\!\!}}
\newcommand\cyan[1]{\colorbox{cyan}{\!\!#1\!\!}}
\newcommand\rmagenta[2][0ex]{\red{\vspace{#1}\magenta{\phantom{:}}\,: #2}}
\newcommand\ryellow[2][0ex]{\red{\vspace{#1}\yellow{\phantom{:}}\,: #2}}
\newcommand\rgreen[2][0ex]{\red{\vspace{#1}\green{\phantom{:}}\,: #2}}
\newcommand\rcyan[2][0ex]{\red{\vspace{#1}\cyan{\phantom{:}}\,: #2}}

%For separated lists with consecutive numbering
\newcounter{separated}

\newcommand{\excise}[1]{}
\newcommand{\comment}[1]{{$\star$\sf\textbf{#1}$\star$}}

%%%%%%%%%%%%%%%%%%%%%%%%%%%%%%%%%%%%%%%%%%%%%%%%%%%
%                                                 %
% TO ENABLE GRADING, DO NOT ALTER ABOVE THIS LINE %
%                                                 %
%%%%%%%%%%%%%%%%%%%%%%%%%%%%%%%%%%%%%%%%%%%%%%%%%%%

%new math symbols taking no arguments
\newcommand\0{\mathbf{0}}
\newcommand\CC{\mathbb{C}}
\newcommand\FF{\mathbb{F}}
\newcommand\NN{\mathbb{N}}
\newcommand\QQ{\mathbb{Q}}
\newcommand\RR{\mathbb{R}}
\newcommand\ZZ{\mathbb{Z}}
\newcommand\bb{\mathbf{b}}
\newcommand\kk{\Bbbk}
\newcommand\mm{\mathfrak{m}}
\newcommand\xx{\mathbf{x}}
\newcommand\yy{\mathbf{y}}
\newcommand\GL{\mathit{GL}}
\newcommand\FL{{\mathcal F}{\ell}_n}
\newcommand\into{\hookrightarrow}
\newcommand\minus{\smallsetminus}
\newcommand\goesto{\rightsquigarrow}

%redefined math symbols taking no arguments
\newcommand\<{\langle}
\renewcommand\>{\rangle}
\renewcommand\iff{\Leftrightarrow}
\renewcommand\implies{\Rightarrow}
\renewcommand\aa{\mathbf{a}}

%to explain about LaTeX commands:
\newcommand\command[1]{\texttt{$\backslash$#1}}

%new math symbols taking arguments
\newcommand\ol[1]{{\overline{#1}}}

%redefined math symbols taking arguments
\renewcommand\mod[1]{\ (\mathrm{mod}\ #1)}

%roman font math operators
\DeclareMathOperator\im{im}

%for easy 2 x 2 matrices
\newcommand\twobytwo[1]{\left[\begin{array}{@{}cc@{}}#1\end{array}\right]}

%for easy column vectors of size 2
\newcommand\tworow[1]{\left[\begin{array}{@{}c@{}}#1\end{array}\right]}


%%%%%%%%%%%%%%%%%%%%%%%%%%%%%%%%%%%%%%%%%%%%%%%%%%%%%%%%%%%%%%%%%%%%%%
\begin{document}%%%%%%%%%%%%%%%%%%%%%%%%%%%%%%%%%%%%%%%%%%%%%%%%%%%%%%
%%%%%%%%%%%%%%%%%%%%%%%%%%%%%%%%%%%%%%%%%%%%%%%%%%%%%%%%%%%%%%%%%%%%%%


\title{\mbox{}\\[-8ex]Math 403 Homework \#3, Spring 2025\\\normalsize
	Instructor: Ezra Miller\\[-2.5ex]}
\author{Solutions by: ...your name...\\[1ex]
Collaborators: ...list those with whom you worked on this assignment...\\
(1 point for each of up to 3 collaborators who also list you)}
\date{Due: 11:59pm Saturday 1 March 2025}

\maketitle

\vspace{-5ex}\collab{0}%

\noindent
\textsc{Reading assignments}\total{}{42}%covers Lectures 13 - 16
\begin{enumerate}[itemsep=-1ex]\setcounter{enumi}{12}
\item%13.
for Thu.~20 February
  \begin{itemize}
  \item{}%
  [Treil, \S 6.3.3--\S 6.3.4] singular value decomposition
  \item{}%
  [Lax, Chap.7: \S Norm of a Linear Map, \S Spectral Radius]
  \item{}%
  [Lax, Chap.8: \S Norm and Eigenvalues]
  \item{}%
  [Treil, \S 6.4] Applications of SVD: spectral radius, operator norm,
  condition number
  \end{itemize}

\item%14.
for Tue.~25 February
  \begin{itemize}
  \item{}%
  [Stewart--Sun, \S II.2.1] matrix norms, consistency
  \item{}%
  [Stewart--Sun, \S IV.1 through Thm 1.3] general perturbation theorems
  \end{itemize}

\item%15.
for Thu.~27 February
  \begin{itemize}
  \item{}%
  [Stewart--Sun, \S IV.1.2 through Thm 1.6] Bauer--Fike theorem
  \item{}%
  [Lax, Appendix 7] Gershgorin's Theorem
  \item{}%
  [Stewart--Sun, \S IV.2 through \S IV.2.1] Gerschgorin theory
  \end{itemize}

\item[16--17.]\addtocounter{enumi}{2}%16. and 17.
for Tue.~4 March and Thu.~6 March
  \begin{itemize}
  \item{}%
  [Tapp, Chap.5] Lie algebras as tangent spaces to the identity
  \item{}%
  [Lax, Chap.9: Matrix-valued Functions, through \S Matrix Exponential]
  \item{}%
  [Wikipedia, \emph{Exponential map (Lie theory)}]
  \end{itemize}
\end{enumerate}

\noindent
\textsc{Exercises}% covers Lectures 9-13 + a bit of 14

\begin{enumerate}

\item\ \score{}%1.
Prove that a subset of a manifold is open if and only if its preimage
under every chart is an open subset of an affine space.  (Note: This
can be used to specify the topology on a manifold without
% appealing to a general notion of topological space
knowing what a topological space is, since affine spaces have the
usual
% topology
notion of ``open'': what it means for a neighborhood of a point in a
manifold to be open is well defined independent of which charts are
used~to~verify~openness.)

\textbf{Solution}

Assume that a subset of the manifold is open but its preimage
is closed. We can also assume that the preimage is not empty, since the 
empty set is open. Then this implies that the image of a closed subset is open
under a homeomorphism, which is false, since homeomorphisms preserve topological strucutre(maps open sets to open sets).
Therefore the preimage of an open sets under every chart is open, likewise same reasoning implies
that the image of an open subset is open.
\item\ \score{1}%2.
Prove that the sphere $S^2 = \big\{\xx \in \RR^3 \,\big|\, ||\xx|| =
1\big\}$ is a smooth manifold.  Is it a rational algebraic variety \red{You need to check inverses. Also differentiable doesn't imply smooth.}
over the field~$\RR$?

\textbf{Solution}

Chart 1:

\begin{align*}
  U_{1}&= \RR^2\\
  \pi_{1}&: U_{1} \hookrightarrow S^2\\
  \pi_{1}^{-1}&: S^2/{(0,0,1)} \twoheadrightarrow U_{1} \simeq R^2 \\
  \pi_{1}^{-1} : (x,y,z)&\mapsto (\frac{x}{1-z},\frac{y}{1-z})\\
  \pi_{1} : (x,y)&\mapsto (\frac{2x}{1+x^2+y^2},\frac{2y}{1+x^2+y^2},\frac{-1+x^2+y^2}{1+x^2+y^2})
\end{align*}

Chart 2:

\begin{align*}
  U_{2}&= \RR^2\\
  \pi_{2}&: U_{2} \hookrightarrow S^2\\
  \pi_{2}^{-1}&: S^2/{(0,0,-1)} \twoheadrightarrow U_{2} \simeq R^2 \\
  \pi_{2}^{-1} : (x,y,z)&\mapsto (\frac{x}{1+z},\frac{y}{1+z})\\
  \pi_{2} : (x,y)&\mapsto (\frac{2x}{1+x^2+y^2},\frac{2y}{1+x^2+y^2},\frac{1-x^2-y^2}{1+x^2+y^2})                
\end{align*}

The transition functions are $\pi_{1} \circ \pi_{2}^{-1}$ 
and $\pi_{2} \circ \pi_{1}^{-1}$, This functions are differeniable, so it is 
a smooth manifold. I don't know what a rational algebraic variety is, but since the 
transition functions are rational I suppose it is a rational algebraic variety.

\item\ \score{}%3.
Fix a field~$F$.  Prove that the unipotent lower-triangular matrices
$U^- \subseteq \GL_n(F)$ map injectively to the set $\FL(F)$ of
complete flags in~$F^n$ expressed as the set of orbits of the group
$B^+$ acting on the right of~$\GL_n(F)$; that is, $U^- \into \FL(F) =
\GL_n(F)/B^+$.

\textbf{Solution}

We need to prove that any lower traingular unipotent matrix corresponds
to one and only one complete Flag. Let's assume the opposite, then you get
two matrices $M_{1}$ and $M_{2}, M_{1}\neq M_{2}$, and the complete flags
induced(Assuming the standdard induciton, where $v_i$ as the span of the union 
of the left $i$ columns) by the two matrices are the same.

Look at the first column of $M_{1}$, since this has to be a linear mulitple
of the first column of $M_{2}$, and since the first entries are the same 
and non-zero(it is 1, because it is lower-triangular and unipotent, therefore
the eigenvalues 1, are on the diagonal), then the first columns are equal.

Assuming that all previous columns are equal, what about the $l^{th}$ column.
Since all the columns are linearly independent in both $M_1$ and $M_2$(they are invertible)
, then you can descirbe $l^{th}$ column as a unique linear combination of the 
$1^{st}$ to $l^{th}$ column. Since the coordinate of $l^{th}$ column in this basis is just one and everything
else being 0 in both matrices, and knowing that the subspace spanned being 
the same, this leads us to conclude that the columns are linear combinations
of each other, and since they both share the same non-zero entry, then they are the 
the same column. Therefore contradicitng that two different unitary matrices
can describe the same complete flag.

\item\ \score{}%4.
Let $S_n \subseteq \GL_n(F)$ be the set of permutation matrices.
Prove that $\FL(F)$ is a rational algebraic variety with atlas
$\{\pi_w: wU^- \to \FL(F) \mid w \in S_n\}$ naturally indexed
by~$S_n$.  You will need to use that any level set
% (actually, you should only need the zero level set)
of a polynomial with coefficients in~$F$ is closed.

\textbf{Solution}

Cliam \#1: Any complete flag can be described as a composition of a lower triangular unipotent matrix 
and permutation matrix.

Since a complete flag can be represented as a sequence of a vectors that are 
continiously added to a subspace, take the first vector to be added or the first 
subspace in the chain. You can scale this vector so that one of its nonzero entries equals 1. Subtract a multiple of the first
vector from the other vectors in the chain so that all entries of the vectors 
corresponding to the 1 in the first vector are 0.

Scale the second vector so that a nonzero entry equals 1, and subtract a multiple of
it, from the other vectors(all vectors other than the first one) so that the entry
corresponding to 1 in the second vector becomes 0.

Continue this process for all vectors. You can always continue this process because
if at any point you find a vector that has all entries 0, then that would imply then 
that vector was in the span of the other vectors, which is a contradiction because
the vectors are all linearly independent.

Put all the current vectors into a matrix $A$, then make a permutation matrix,
where the $i^{th}$ column has a one in the $j^{th}$ place if $c_{i}$(the $i^{th}$ column of A), had 1 chosen at the $j^{th}$ entry.

When you multiply A with this permutation matrix you get a lower traingular unipotent matrix,
since the $1^{s}$ are going to be on the diagonal and anything after that entry will become 0
from the construction. The inverse of the permutation matrix created will be our chart.

Claim \#2: The transition functions between charts are composition of permutation matrices.

To get the transition function between the charts compsose the matrix with the inverse of the 
other matrix, which is just another permutation matrix.

Claim \#3: This manifold is a rational algebraic variety.

I don't what rational algebraic varities are, but since you the composition of the 
matrices is rational, I suppose it is a rational algebraic variety.
\item\ \score{}%5.
The \emph{standard complex structure} on $\RR^{2n}$ is the
block-diagonal $2n \times 2n$ matrix $J_{2n}$ whose diagonal blocks
are all $\twobytwo{0&\!-1\\1&\ 0}$.  Prove that $A \in \GL_{2n}(\RR)$
is complex-linear if and only if $A$ commutes with the standard
complex structure: $AJ_{2n} = J_{2n}A$.

\textbf{Solution}

For $n=1$, assume there exists some matrix A and it is commutative with the standard complex strucutre. Let's what
conditions it imposes.

$$
A = \begin{bmatrix}
  a & c\\
  d & b
\end{bmatrix}
$$

$$
\begin{bmatrix}
  0 & -1\\
  1 & 0
\end{bmatrix}
\begin{bmatrix}
  a & c\\
  d & b
\end{bmatrix}=
\begin{bmatrix}
  a & c\\
  d & b
\end{bmatrix}
\begin{bmatrix}
  0 & -1\\
  1 & 0
\end{bmatrix}
$$

$$
\begin{bmatrix}
  -d & -b\\
  a & c
\end{bmatrix}=
\begin{bmatrix}
  c & -a\\
  b & -d
\end{bmatrix}
$$
So this means commutatitve implies,\\
\begin{align*}
a=b\\
-d=c 
\end{align*}

Which is the same coniditions for a 2$\times$2 to be complex linear.
Note this proves both sides of the iff since, commutativity and complex linearity impose the same restrictions.

Assume true for size n, now we will prove that it is true for n+1. We only have
to check for complex linearity in each 2$\times$2 block in the matrix. So if 
each 2$\times$2 commutes with the standard complex structure, then the entire 
matrix is complex linear.

\begin{align*}
  \begin{bmatrix}
    0 & -1 & ...\\
    1 & 0 & ...\\
    ... & ...& J_{2n}\\
  \end{bmatrix} A&= 
  A\begin{bmatrix}
    0 & -1 & ...\\
    1 & 0 & ...\\
    ... & ...& J_{2n}\\
  \end{bmatrix} \\
  \begin{bmatrix}
    0& -1& ... 0\\
    1& 0& ... 0\\
    ... 0& 0&... J_{2n}
  \end{bmatrix}
  \begin{bmatrix}
    A_{11} & A_{12}& ... A_{1n}\\
    A_{21} & A_{n}\\
    \dots\\
    A_{1n}
  \end{bmatrix}&=
  \begin{bmatrix}
    A_{11} & A_{12}& ... A_{1n}\\
    A_{21} & A_{n}\\
    \dots\\
    A_{1n}
  \end{bmatrix}\begin{bmatrix}
    0& -1& ... 0\\
    1& 0& ... 0\\
    ... 0& 0&... J_{2n}
  \end{bmatrix}\\
  \begin{bmatrix}
    A^{\vert}_{11} & A^{\vert}_{12}& ... A^{\vert}_{1n}\\
    A^{\vert}_{21} & A^{\vert}_{n}\\
    \dots\\
    A^{\vert}_{1n}
  \end{bmatrix}&=
  \begin{bmatrix}
    A^{-}_{11} & A^{-}_{12}& ... A^{-}_{1n}\\
    A^{-}_{21} & A^{-}_{n}\\
    \dots\\
    A^{-}_{1n}
  \end{bmatrix}
\end{align*}  

Where $A^{\vert}=A^{-}$, implies that the 2$\times$2 matrix is complex linear(from the $n=1$ case).
We now see that the matrix is commutatitve if and only if every 2$\times$2 is complex linear
which iff the matrix is complex linear.


\item\ \score{}%6.
Prove that if $\lambda \in \CC$ is an eigenvalue of a unitary matrix
then $|\lambda| = 1$.

\textbf{Solution}

\begin{align*}
  U^{*}U=UU^{*}=I\\
\end{align*}

Let $x \in $ eigen-space($\lambda$).
\begin{align*}
  Ux=\lambda x \rightarrow x^{*}U^{*}=x^{*} \overline{\lambda}\\
  x^{*} U^{*}Ux=(x \overline{\lambda})\lambda x\\
  x^{*}Ix=\vert\lambda| x^{*}x\\
  |x|= |\lambda||x|, x\neq0\\
  1=|\lambda|
\end{align*}

\item\ \score{}%7.
Show that the standard Hermitian inner product on $\CC^n$ defines a
distance $d(\xx,\yy) = ||\xx-\yy||$ on~$\CC^n$.  Give an example of an
isometry $\varphi$ of~$\CC^n$ such that $\varphi(\0) = \0$ but
$\varphi$ is not $\CC$-linear.

Properties of distance\\
1) d(x,x)=0, already satisfied\\
2) if $x\neq y$, then d(x,y) $>$ 0. Since only 0 has absolute
value 0, this property is satisfied.
3) d(x,y)=d(y,x)\\
$||x-y||$=$\sqrt{(x-y) \cdot (x-y)} = \sqrt{(y-x) \cdot (y-x) } = ||y-x||$\\

4) Triangle Property
$d(x,y)\leq d(x,z) + d(z,y)$
Since norms also satisfy the triangle property $|u+v| \leq |u| + |v|$.
Then let $|u+v|=|x-y| \wedge |u|=|x-z| \wedge |v|=|z-y|$. So the triangle property is fulfilled.


\item\ \score{}%8.
Let $\varphi$ be an orthogonal transformation of a real inner product
space~$V$.  Assume that $\varphi^2 = -I$.  Show that $\dim V$ is even,
say~$2n$.  Moreover, prove that there exists a dimension~$n$ subspace
$W \subset V$ and an
% orthogonal transformation
isometry $\psi: W \to W^\perp$ such that, in the decomposition $V = W
\oplus W^\perp$, the operator $\varphi$ is given by the block matrix
$$%$$ $
  \twobytwo{ \0  & -\psi^*
          \\ \psi &    \0 }
$$%$$ $
(N.B. The result means that $\varphi$ can be thought of as
multiplication by~$i$ on a complex vector space whose real and
imaginary parts are $W$ and~$W^\perp$.)

\textbf{Solution}

If dimension of V is odd, then after applying the transformation twice,
the basis can't be left-handed because either the transformation
changes signs twice or not at all, and in either case the matrix has to have positive
determinant after applying the transformation twice, so the dimension of V can't be odd.

Now let's prove that the specified subspace exists in the ambient space.

Claim 1: $\varphi (x) \perp x$.

\begin{align*}
  \varphi^2(x)&=-x,\varphi(-\varphi) = I\\
  \varphi^{-1}&=-\varphi=\varphi^{T}\\
  x^{T}(\varphi^{2})^T&=(\varphi^2 x)^T = (-x)^T\\
  x\cdot&\varphi x\\
  -\varphi^{2}x&\cdot \varphi x\\
  x^{T} \varphi^{T}&\cdot -\varphi^{2} x\\
  x^{T} -\varphi&\cdot -\varphi^{2} x\\
  x^{T}\varphi^{3}x&\\
  -x^{T}\varphi x=&x \cdot \varphi x \rightarrow x\varphi x=0
\end{align*}

The number can only be equal to both its positive and negative iff the number 
is 0.

Claim 2: We can construct W.\\
Pick $V_1$ and $\varphi(V_1)=V_1^{'}$.\\
Pick $V_2 \perp V_1, V_1^{'}$, then $ \varphi(V_2)\perp V_1$ and $V_1^{'}$, since
 the inner product is conserved under orthogonal transformations. Keep making 
pairts until you have n pairs, where 2n is the dimension of V. Now take $W$ to be 
the subspace defined by taking the span of $V_1, V_2, ..., V_n$. Then $\varphi$ takes
$W$ to $W^{\perp}$, and induces $\psi$, which is just a restriction of the domain 
to $W$ and restrction of the codomain to $W^{\perp}$. Since the inner product is conserved in $W^{\perp}$,
the distance between points is also conserved and therefore $\psi$ is also an isometry.

Any $x \in V$ can be decomposed into an element in $W$ and $W^{\perp}$. 
Therefore we just need to figure out where $\psi$ takes the elments from the $W$.
Let's take a basis $v \in W$, and its corresponding pair $v^{'}$.

\begin{align*}
  \psi (v)&= v^{'} \text{\; By Defintion.}\\ 
  \psi^{*}(v^{'})&=v \text{\; Because }\varphi^{*}=\varphi^{-1}
\end{align*}

So make the block matrix 
$$%$$ $
  \twobytwo{ \0  & -\psi^*
          \\ \psi &    \0 }
$$%$$ $

Which takes $x=w+w^{'}$ to $w' - w$, and $w' - w$ to $ -w' - w$, which satisifies
the original function.

\item\ \score{}%9.
True or false: the sum of two normal operators is normal.  Justify.

\textbf{Solution}

Take this counterexample. Where A and B, are normal opertors.   

$$
A=
\begin{bmatrix}
  5& 0\\
  0& 9
\end{bmatrix}
$$
$$
B=
\begin{bmatrix}
  0& -3\\
  3& 0
\end{bmatrix}\\
$$
$$
A+B=
\begin{bmatrix}
5&-3\\
3&9
\end{bmatrix}
$$
\begin{align*}
  (A+B)(A+B)^{*}=
  \begin{bmatrix}
    34& -12\\
    -12& 90
  \end{bmatrix}\\
  (A+B)^{*}(A+B)=
  \begin{bmatrix}
  34& 12\\
  12& 90
\end{bmatrix}    \\
 (A+B)(A+B)^{*}\neq(A+B)^{*}(A+B)\\
\end{align*}

So the statement is not true.
\item\ \score{}%10.
Show that the space of positive (semi)definite real symmetric
matrices is convex.  Is the same true with ``complex Hermitian'' in
place of ``real symmetric''?

\textbf{Solution}

Assume A and B are postive semi definite matrices. Then A and B, satisfy

$$x^T A x \geq 0 \; \forall x$$

Now we want to show that $\alpha$A+$(1-\alpha)$B, also satisfies this property.

$$x^{T}(\alpha A+ (1-\alpha) B)x = \alpha x^T Ax + (1-\alpha) x^T Bx \geq 0 \; \forall x \wedge 0 \leq \alpha \leq 1$$

So the real positive matrices are convex. This proof also works for conjugate transposes too.



\item\ \score{}%11.
Orthogonally diagonalize the matrix $A = \twobytwo{3&2\\2&3}$.  Find
all square roots of~$A$; note that they are all self-adjoint.

\textbf{Solution}

Since the trace of the matrix is 6, $\lambda_1 + \lambda_2 = 6$. If we let
$\lambda=1$, $(A-I)$, becomes singular. So the other eigen value becomes 5.

The eigen vector in eigen space of 1, is (1,1), and the other eigen vector is (-1,1).

$$
P=
\begin{bmatrix}
  1&-1\\
  1& 1
\end{bmatrix}
$$
$$
P^{-1}=
\begin{bmatrix}
  -1/2& 1/2\\
  1/2& 1/2
\end{bmatrix}
$$
$$
D=
\begin{bmatrix}
  5& 0\\
  0& 1
\end{bmatrix}
$$

$$
\sqrt{D}=
\begin{bmatrix}
  \sqrt{5}&0\\
  0& 1
\end{bmatrix}
$$

$$
\sqrt{A}=P\sqrt{D}P^{-1}
$$

\item\ \score{}%12.
Find a singular decomposition of the matrix $A = \twobytwo{2&3\\0&2}$.
Use it to find $\max_{\Vert\xx\Vert \leq 1} \Vert A\xx\Vert$ and
$\min_{\Vert\xx\Vert = 1} \Vert A\xx\Vert$, as well as the vectors
where this maximum and minimum are attained.  Describe geometrically
the image under $A$ of the closed unit disk in~$\RR^2$.

\textbf{Solution}

$$
A^{T}A=
\begin{bmatrix}
 2& 0\\
 3& 2 
\end{bmatrix}
\begin{bmatrix}
  2&3\\
  0&2\\
\end{bmatrix}
=
\begin{bmatrix}
  4&6\\
  6&13
\end{bmatrix}
$$

The eigenvalues of are 16 and 1, therefore the singular values are 
4 and 1. Therefore the norm operator is 4, and the minimum length of unit vectors
is 1.

The eigen valeus of $A^{T}A$ are $(1,2)$ and $(-2,1)$ for the singular values 4 and 1 
respectively. So the vector that achieve the maximum and the minimum are precisely the normalized version of 
this vectors.

From this we also find $V^{*}$.

$$
V=\begin{bmatrix}
  1/\sqrt5 & -2/\sqrt5\\
  2/\sqrt5& 1/\sqrt{5}
\end{bmatrix}
$$

To W we just do $w_[i]=\frac{1}{\sigma_{i}}A v_{i}$.

$$
W=\begin{bmatrix}
  2/\sqrt5 & -1/\sqrt5\\
  -1/\sqrt5 & 2/\sqrt5
\end{bmatrix}
$$
So $A=W \Sigma V^{*}$, where $\Sigma$ is a just an operator is the singular values in the diagonal.

\item\ \score{}%13.
Prove that the operator norm of a matrix~$A$ coincides with the
Frobenius norm of~$A$ if and only if $A$ has rank at most~$1$.

\textbf{Solution}

If A has rank 0, there is nothing to prove. The frobenius norm of a matrix is
the largest singualr value of the matrix(From the notes). If A has rank one, that
means it only has 1, non-zero singular value, therefore that value has to be the 
one that has to outputed by the frobenius norm.

If the matrix has rank greater than 1, then the frobenius norm and the 
operator norm can't be equal since the frobenius norm is the sqaure root of the sum 
of the squares of the singular values, which is greater than
the heighest singular value. There the relationship is iff.

\end{enumerate}


%%%%%%%%%%%%%%%%%%%%%%%%%%%%%%%%%%%%%%%%%%%%%%%%%%%%%%%%%%%%%%%%%%%%%%
\end{document}%%%%%%%%%%%%%%%%%%%%%%%%%%%%%%%%%%%%%%%%%%%%%%%%%%%%%%%%
%%%%%%%%%%%%%%%%%%%%%%%%%%%%%%%%%%%%%%%%%%%%%%%%%%%%%%%%%%%%%%%%%%%%%%


%%% Various Emacs customizations:
%%% Local Variables:
%%% fill-column: 70
%%% indent-tabs-mode: t
%%% TeX-electric-sub-and-superscript: nil
%%% TeX-brace-indent-level: 0
%%% LaTeX-indent-level: 0
%%% LaTeX-item-indent: 0
%%% TeX-master: t
%%% End: